% =====================================================================
% 参考文献
% ---------------------------------------------------------------------
% 说明：
%   1. 本文正文采用"作者（年份）"的行文方式，未使用 \cite，故此处用
%      thebibliography 直接列出，编译后即为参考文献表。
%   2. 全部条目的出处与卷期已核验，核验过程见 docs/文献引用核验表.md。
%   3. 中文条目按第一作者拼音排序，外文在前。
% =====================================================================

\begin{thebibliography}{99}
\footnotesize

\bibitem{lucca2009}
Lucca D O, Trebbi F. Measuring Central Bank Communication: An Automated
Approach with Application to FOMC Statements. NBER Working Paper No.~15367, 2009.

\bibitem{shapiro2022}
Shapiro A H, Wilson D J. Taking the Fed at Its Word: A New Approach to
Estimating Central Bank Objectives Using Text Analysis.
\textit{Review of Economic Studies}, 2022, 89(5): 2768--2805.

\bibitem{chen2024}
陈妍玲, 唐靖, 郭文伟. 央行沟通与金融系统性风险：促进还是抑制？
\textit{上海立信会计金融学院学报}, 2024, 36(6): 3--20.

\bibitem{ji2026}
纪钧尧, 尤晓彤, 陈旭. 央行货币政策文本情绪对汇率的影响：基于 FinBERT
模型的多国文本情绪挖掘. \textit{世界经济研究}, 2026(7).

\bibitem{jiang2021}
姜富伟, 胡逸驰, 黄楠. 央行货币政策报告文本信息、宏观经济与股票市场.
\textit{金融研究}, 2021(6): 95--113.

\bibitem{jiang2020}
姜富伟, 孟令超, 唐国豪. 媒体文本情绪与股票回报预测.
\textit{经济学（季刊）}, 2020, 已录用待刊.

\bibitem{jiang2019}
Jiang F, Lee J, Martin X, Zhou G. Manager Sentiment and Stock Returns.
\textit{Journal of Financial Economics}, 2019, 132(1): 126--149.

\bibitem{jiang2025}
蒋海, 孙娜, 李欣明. 央行预期管理能缓解银行流动性囤积吗？——基于前瞻性
文本情绪指数的实证检验. \textit{金融研究}, 2025(6): 1--20.

\bibitem{liao2026}
廖海龙. 大模型驱动的情绪因子构建及其在 A 股量化择时中的应用.
江西财经大学硕士学位论文, 2026.

\bibitem{lin2025}
林建浩, 孙乐轩. 大语言模型与经济金融文本分析：基本原理、应用场景与研究展望.
\textit{计量经济学报}, 2025, 5(1).

\bibitem{wang2023}
王辉, 朱家雲, 胡诣聪. 央行预期引导可以降低银行系统性金融风险吗？——基于
市场解读偏离的视角. \textit{金融研究}, 2023(9): 1--19.

\bibitem{wang2021}
王熙, 郑梦圆. 中国货币政策冲击的信息效应与“价格之谜”——基于文本分析的
研究. \textit{山东大学学报（哲学社会科学版）}, 2021(5).

\bibitem{wang2025}
王宇伟, 刘宏雅, 范从来. 央行言辞沟通能否提振企业家信心.
\textit{经济理论与经济管理}, 2025, 45(8): 33--52.

\bibitem{wang2022}
王春芝. 货币政策公告的股价预漂移效应研究——基于风险敏感理论的视角.
南京财经大学硕士学位论文, 2022.

\bibitem{wangyu2025}
汪禹, 吴彬. 情感分析视角下中文政策文本的倾向词典构建研究.
\textit{情报资料工作}, 2025, 46(5).

\bibitem{xu2023}
许魁. 央行货币政策报告文本信息与股票市场——基于分析师报告视角.
山东大学硕士学位论文, 2023.

\bibitem{you2026}
游庭玮. 美联储政策利率对 A 股行业收益的冲击效应研究.
华东政法大学硕士学位论文, 2026.

\bibitem{zhou2024}
周阳, 綦岚鑫. 央行文本情绪具有金融市场效应吗——基于《中国货币政策执行
报告》的研究. \textit{财经科学}, 2024(2).

\bibitem{zhou2018}
周亮. 货币周期对我国股市行业轮动的影响.
\textit{湖北经济学院学报}, 2018(3): 31--37.

\end{thebibliography}
